\section{One residual seed for all later towers}

% BEGIN S66C1
\paragraph{S66C1. Discharging the arithmetic inputs on the original tower.}
Fix the original Case III datum $D$, a tower $T_0$ satisfying
$\mathrm{Ok}$ and $\mathrm{Four}$, the comparison $\iota$, and
the ambient reduction datum $\mathfrak a$. Use the imaginary
quadratic field $K'$ specified by $\mathrm{Four}$ and the original
mark of $T_0$. The mark is outside $G_{K'}$ by the two distinct
diagonal characters at the mark (S65A1). The Case III hypotheses give $\ell\ne2$ and the specified rational
complex conjugation. The quadratic-field datum gives degree two, total
imaginarity, $\ell\nmid\mathrm{disc}(K')$, and splitting at $\ell$.
Together with the mark, these give all the inputs of D66T8. That
construction needs
the global reciprocity map proved in CIP1. It gives the positive
integer $m$, one coefficient ring $O_0$, and the promise for
\emph{every} later coefficient ring and positive power of $\kappa$.
This is the only application of the standing
$\mathrm{Af.Hypotheses}'$ in the source assembly: its sole occurrence
after the hypotheses are introduced is as an argument to the
dihedral constructor. We use the explicit constructor D66T8 here.

P66E5 supplies $\mathrm{PropEGP}$. Apply R66P6 with the prescribed
coefficient lower bound $O_0$ and numerical lower bound $18m+3$.
This fixes $w,\theta,k,\rho$, the integral lattice, all its specified
local slopes, and a ring $O_1$ extending both coefficient rings.
The $\mathrm{PR}$ property of this $\rho$ and its irreducibility
on every open subgroup are the hypotheses for F66A1. Realize
$\kappa$ at power one over $O_1$ first. F66A1--F66A4 then fix one
$j\in\{1,\ldots,9\}$ which works for every integral realization
at that power. Return to D66T8 at power $j$ over the same $O_1$.
For the resulting tensor the attachment and the all-place window
of F66A4 hold together. Since $2mj+3\le18m+3\le k$, both outside
gaps and the middle gap in the four weights satisfy S65A2.
The remaining twist, the auxiliary prime, and the reduction map are
those of S65A1. Every choice thus far precedes the variable tower $T$.
% END S66C1

% BEGIN S66C2
\paragraph{S66C2. The original Case III seed.}
Let $Q_T$ denote the integral signed tensor with its scalar twist
on the tower $T$, as in S65A2--S65A3. Choose
\[
              s = \mathfrak a.\mathrm{red}_{O_1}(Q_{T_0})
\]
entry by entry on the original group. The joint input of S66C1
makes the conclusions of S65A3 available on that exact tensor and
ring. On $T_0$ the residual pair has allowed auxiliary set $\{w\}$,
the required diagonal characters, continuity, and the stipulated
$N_M$ support. For every later acceptable $T\ge T_0$ satisfying
$\mathrm{Dies}(\ell,w,T.F)$ the character $\theta$ is trivial on
$G_{T.F}$. Thus $Q_T$ satisfies the integral pair with empty
auxiliary set and the weight and automorphic clauses of the seed.
Its reduction is compatible with the prechosen $s$ on the indicated
overlap, has type $A$, and realizes both required nonzero root
entries. These are separate assertions in the source, all supplied
by S65A3. In particular the existential witnesses $w,s$ are
outside the quantifier on $T$. This proves the full
$\mathrm{SeedIVAlt}$ contract in the human reconstruction using
the stated classical arithmetic theorems.
% END S66C2
